% chapters/03_learned_representations.tex
%
% Part I, Chapter 3.  Learned representations: PCA, k-means, matching
% pursuit, neural networks (FC and conv).  Adapted from the geometric
% tour's Chapter 3 with matching pursuit added (it's central to Part III)
% and the "survey of architectures" section cut (deferred to Part II,
% which builds those architectures in depth, and the closing chapter of
% Part III, which explains what we deliberately didn't cover).

\chapter{Learned representations}
\labch{learned}

\epigraph{All models are wrong, but some are useful.}{George E.~P.~Box}

%------------------------------------------------------------------------
\section{When the basis comes from data}
%------------------------------------------------------------------------

\noindent The previous chapter built three representations whose
reference vectors --- sinusoids for Fourier, scaled and shifted wavelets
for the DWT, products of those for Scattering --- were chosen by
mathematicians in advance.  This works beautifully when we know what
property we want to expose.  We want pitch, so we project onto
sinusoids.  We want time-frequency structure, so we project onto
wavelets.  The chosen reference vectors carry our prior knowledge about
the signal.

\marginnote{%
\textbf{The bias-variance trade-off.}  Predefined bases are
\emph{low-variance}: they don't depend on which data you fed them, so
they don't overfit.  They're also \emph{high-bias}: if your signal
doesn't decompose well into the chosen basis, no amount of clever
analysis will help.  Learned bases reverse the trade-off.  This is the
classical bias-variance picture, made concrete.}

But what do we do when we don't know what to look for?  When we have a
million labeled images of cats and dogs and want a representation that
distinguishes them, we have no idea what mathematical function picks
out the relevant cat-vs-dog features.  When we have orchestral recordings
and want a representation that captures a composer's style, we cannot
write down by hand what reference vectors would expose ``Mahler-ness.''
The features are real but they are not in any catalogue we know how to
consult.

The natural move is to \emph{let the data choose the basis}.  Rather
than designing the reference vectors mathematically, we look at many
examples of signals and derive a set of vectors that, in some sense,
captures what is common or what is distinguishing.  This is what
\emph{learned representations} do.  The basis is no longer in a textbook;
it lives in a parameter array computed from data.

This chapter introduces four learned representations of increasing
sophistication.  Principal Component Analysis (Section~\ref{sec:pca})
finds the directions of greatest variance in a dataset.  $K$-means
(Section~\ref{sec:kmeans}) finds prototypes that group similar
examples.  Matching pursuit (Section~\ref{sec:mp}) builds a sparse,
interpretable approximation of a signal from a large dictionary of
candidates.  Neural networks (Section~\ref{sec:nn}) learn nonlinear,
multi-layered representations whose components are themselves
learned-from-data projections.

In all four cases, the underlying structure is the same as in the
previous chapter: project the signal onto reference vectors, possibly
apply a nonlinearity, possibly stack layers.  The only difference is
where the reference vectors come from.

%------------------------------------------------------------------------
\section{Principal Component Analysis}
\label{sec:pca}
%------------------------------------------------------------------------

PCA is the simplest learned representation: a linear projection whose
basis vectors are the directions of maximum variance in the data.

Given a data matrix $\mathbf{X} \in \R^{N \times d}$ with one example per
row (so $N$ examples of dimension $d$), PCA proceeds in three steps.

\paragraph{1. Centre the data.}  Subtract the per-feature mean so that
the columns of $\mathbf{X}$ have zero mean.  This is essential because
the geometric machinery of PCA assumes the data is centred at the
origin.

\paragraph{2. Compute the covariance.}  Form
\begin{equation}
    \mathbf{C} \;=\; \frac{1}{N}\, \mathbf{X}^\top \mathbf{X}.
\end{equation}
The $(i, j)$ entry of $\mathbf{C}$ is the covariance between feature $i$
and feature $j$.

\paragraph{3. Diagonalise.}  Find the eigenvalues $\lambda_1 \ge
\lambda_2 \ge \cdots \ge \lambda_d \ge 0$ and corresponding eigenvectors
$\mathbf{e}_1, \ldots, \mathbf{e}_d$ of $\mathbf{C}$.  These eigenvectors
are the \emph{principal components} --- the directions in feature space
along which the data has the most variance.  The first eigenvector
$\mathbf{e}_1$ is the single direction along which the data spreads
most; $\mathbf{e}_2$ is the direction of next greatest spread
perpendicular to $\mathbf{e}_1$; and so on.

To get a $k$-dimensional representation we project onto the first $k$
eigenvectors.  Let $\mathbf{E}_k \in \R^{d \times k}$ be the matrix
whose columns are $\mathbf{e}_1, \ldots, \mathbf{e}_k$.  Then
\begin{equation}
    \mathbf{Z} \;=\; \mathbf{X}\, \mathbf{E}_k
\end{equation}
is the data expressed in the $k$-dimensional principal-component basis.

\marginnote{%
PCA is mathematically a singular value decomposition (SVD) of the data
matrix $\mathbf{X}$ in disguise: the principal components are the right
singular vectors of $\mathbf{X}$.  When we work with PCA in Part~II,
the SVD is the practical computation.}

\paragraph{PCA in the framework.}  Compare $\mathbf{Z} = \mathbf{X}\,
\mathbf{E}_k$ with the analysis formula
$\mathbf{c} = \mathbf{B}\, \xb$ from the previous chapter.  PCA is
exactly that, with $\mathbf{B} = \mathbf{E}_k^\top$ (the principal-component
matrix) and the data being analysed simultaneously rather than one
example at a time.  The reference vectors are the eigenvectors.  The
only thing that distinguishes PCA from the Fourier transform, at the
level of mathematical structure, is that the basis vectors come from
the data instead of from $e^{i 2 \pi k n / N}$.

\paragraph{What PCA captures.}  PCA is a tool for finding low-dimensional
linear structure in high-dimensional data.  When applied to natural
images, its principal components often look like blurry blobs and
gradients --- the most-variance directions in the space of $32 \times 32$
RGB images.  When applied to speech features, the principal components
capture coarse vowel-vs-consonant structure.  When applied to text
embeddings, they capture broad semantic axes (topic, sentiment).

What PCA cannot capture is nonlinear structure.  The data may lie on
a curved manifold in feature space; PCA will fit it with a flat
hyperplane and lose all the curvature.  For nonlinear structure we
need either to engineer nonlinear features by hand (a hard problem) or
to learn them with a network (the rest of this chapter).

%------------------------------------------------------------------------
\section{$K$-means and prototypes}
\label{sec:kmeans}
%------------------------------------------------------------------------

$K$-means takes a different geometric view.  Rather than finding
directions of variance, it finds \emph{prototypes}: a small set of
representative points such that every other data point is close to one
of them.

\paragraph{The algorithm.}  Given data points $\xb_1, \ldots, \xb_N$
and a chosen number of clusters $K$:

\begin{enumerate}
    \item Initialise $K$ prototype vectors $\boldsymbol\mu_1, \ldots,
    \boldsymbol\mu_K$ (often by randomly picking $K$ of the data
    points).

    \item \emph{Assign} each data point $\xb_n$ to the nearest
    prototype:
    \[
        C_n \;=\; \arg\min_k \norm{\xb_n - \boldsymbol\mu_k}.
    \]

    \item \emph{Update} each prototype to the mean of the data points
    assigned to it:
    \[
        \boldsymbol\mu_k \;=\; \frac{1}{|\{n : C_n = k\}|}
        \sum_{n : C_n = k} \xb_n.
    \]

    \item Repeat steps 2--3 until the assignments stop changing.
\end{enumerate}

\paragraph{$K$-means in the framework.}  $K$-means looks like a
different kind of object from PCA --- a discrete clustering rather than
a linear projection.  But viewed through our framework it is still a
representation by reference vectors: every signal $\xb$ is summarised
by the index $C(\xb)$ of its nearest prototype, and the prototypes
$\{\boldsymbol\mu_k\}$ are the reference vectors.  The ``projection''
is the assignment operation, which is \emph{not} a linear map but a
piecewise-constant function of $\xb$.  This is our first nonlinear
representation, but the nonlinearity is hidden in the
\emph{arg-min}.

\marginnote{%
$K$-means and PCA are sometimes presented as alternatives, but they
answer different questions.  PCA: ``in what direction does the data
spread the most?''  $K$-means: ``what are the typical examples?''
A composer's archive analysed with PCA might give you a few axes of
stylistic variation.  Analysed with $K$-means, it might give you a
handful of representative pieces.  Both are useful; the right tool
depends on what you want to do.}

A useful matrix formulation: let $\mathbf{M} \in \R^{K \times d}$ be the
matrix whose rows are the prototypes, and let $\mathbf{C} \in \R^{N
\times K}$ be the binary assignment matrix with $C_{nk} = 1$ if example
$n$ is in cluster $k$ and zero otherwise.  Then the reconstruction
$\hat{\mathbf{X}} = \mathbf{C}\, \mathbf{M}$ replaces each data point
with its assigned prototype, and $K$-means minimises the reconstruction
error $\norm{\mathbf{X} - \mathbf{C}\, \mathbf{M}}_F^2$.  This is
structurally an extreme dimensionality reduction --- each example is
reduced to one integer in $\{1, \ldots, K\}$.

%------------------------------------------------------------------------
\section{Matching pursuit and sparse dictionaries}
\label{sec:mp}
%------------------------------------------------------------------------

The two representations so far --- PCA and $K$-means --- both insist on
a small basis: $k$ eigenvectors or $K$ prototypes.  Matching pursuit
takes the opposite tack: provide a \emph{large} dictionary of candidate
reference vectors, and let the analysis algorithm choose a small subset
of them per signal.  This is the right idea when no fixed-size basis
would do justice to the variety of signals we want to represent.

\paragraph{Setup.}  Let $\mathbf{D} \in \R^{d \times M}$ be a
\emph{dictionary}: a matrix whose $M$ columns $\mathbf{d}_1, \ldots,
\mathbf{d}_M$ are candidate reference vectors (atoms).  We allow $M$
to be much larger than $d$; the dictionary is \emph{overcomplete}.
Atoms could be Gabor wavelets at many scales and frequencies, samples
from a corpus of audio recordings, or anything else relevant.

The goal is to approximate a signal $\xb$ as a sparse linear combination
of atoms:
\begin{equation}
    \xb \;\approx\; \sum_{k=1}^{K} \alpha_k\, \mathbf{d}_{i_k},
    \qquad
    K \ll M.
\end{equation}
We want to find which $K$ atoms (the indices $i_1, \ldots, i_K$) and
which coefficients $\alpha_1, \ldots, \alpha_K$ minimise the
approximation error.  Choosing the best subset of $K$ atoms out of $M$
is combinatorially explosive --- there are $\binom{M}{K}$ subsets to
consider.  In practice we use a \emph{greedy} algorithm.

\paragraph{The matching pursuit algorithm.}  At each step we pick the
atom that explains the largest fraction of what remains.

\begin{enumerate}
    \item Start with residual $\mathbf{r}_0 = \xb$.
    \item For $k = 1, 2, \ldots, K$:
    \begin{enumerate}
        \item Find the atom most correlated with the current residual:
        \[
            i_k \;=\; \arg\max_i \; \abs{\inner{\mathbf{r}_{k-1}}{\mathbf{d}_i}}.
        \]
        \item Compute its coefficient: $\alpha_k = \inner{\mathbf{r}_{k-1}}{\mathbf{d}_{i_k}}$
        (assuming dictionary atoms are unit-norm).
        \item Subtract the projection: $\mathbf{r}_k = \mathbf{r}_{k-1}
        - \alpha_k \mathbf{d}_{i_k}$.
    \end{enumerate}
    \item Stop when $\norm{\mathbf{r}_K}$ is small enough, or when $K$
    reaches a budget.
\end{enumerate}

\noindent The algorithm produces an approximation $\hat\xb = \sum_k
\alpha_k \mathbf{d}_{i_k}$ that uses only $K$ of the $M$ available
atoms.

\paragraph{Why this is a representation.}  Matching pursuit produces,
for each input signal $\xb$, a sparse list of (atom, coefficient) pairs.
This list \emph{is} the representation.  Unlike a dense Fourier
spectrum or a flat PCA projection, matching pursuit returns only
those components that are large; the rest are implicit zeros.  The
representation is naturally \emph{interpretable}: each atom is a
specific pattern, and the coefficient tells you how much of that
pattern is in the signal.

\marginnote{%
Matching pursuit was introduced by Mallat and Zhang in 1993.  Its
descendants include orthogonal matching pursuit (OMP), basis pursuit
(which solves a related $L^1$ minimisation problem), and the broader
field of sparse coding and compressed sensing.  When the dictionary
itself is also learned from data, the algorithm is called
\emph{dictionary learning} or \emph{sparse coding}; this borders on
unsupervised neural-network methods.}

\paragraph{Matching pursuit in the framework.}  Matching pursuit is a
representation by reference vectors --- the atoms --- but the
``projection'' is no longer a single matrix multiplication.  It is the
output of a small algorithm (the greedy selection loop).  The basis is
also no longer a basis in the strict sense: the dictionary is redundant,
and we deliberately use only a few atoms per signal.

This breaks the clean ``matrix in, matrix out'' picture of the previous
chapter, but it gains two things in exchange.  First, \emph{sparsity}:
the representation has very few non-zero components, which is exactly
what we often want for compression and for interpretation.  Second,
\emph{adaptivity}: each signal gets its own custom subset of atoms,
matched to its content, rather than being squeezed through the same
projection as every other signal.

\paragraph{Why matching pursuit matters here.}  We will meet matching
pursuit again in Part~III, where it plays the role of a learned-but-not-neural
codec: a way to decompose orchestral sounds into a sequence of
atom-indices that captures their timbre and structure.  These
atom-indices then become the input to a transformer model that learns
to generate music.  In this sense matching pursuit gives us the best
of both worlds: an interpretable codec like Fourier, and a learned
basis like a neural network.

This pattern --- combine an a-priori transform with a learned model on
its output --- is increasingly common.  It is the basis of AudioLM,
MusicLM, and several other recent music-generation systems.  We will
build our own version in Part~III.

%------------------------------------------------------------------------
\section{Neural networks}
\label{sec:nn}
%------------------------------------------------------------------------

We come to the modern centre of mass.  A \emph{neural network} is a
representation built from many layers, each of which is a learned
linear projection followed by a nonlinearity.

\paragraph{A single layer.}  The basic building block, repeated many
times, is
\begin{equation}
    \zb \;=\; f(\Wb \xb + \mathbf{b}).
    \label{eq:nn_layer}
\end{equation}
Here $\Wb$ is a learned matrix of weights, $\mathbf{b}$ is a learned
bias vector that allows the projection to be affine rather than purely
linear, and $f$ is a fixed nonlinearity applied to each component (most
commonly ReLU, $f(x) = \max(0, x)$).

The inner product structure is right there in $\Wb \xb$.  Each row of
$\Wb$ is a reference vector $\wb_i$, and the $i$-th component of the
output is
\begin{equation}
    z_i \;=\; f\!\left( \inner{\wb_i}{\xb} + b_i \right).
\end{equation}
This is the inner-product analysis of the previous chapter, with two
modifications: the basis vectors $\wb_i$ are learned, and a nonlinearity
$f$ is applied to each result.

\paragraph{Logistic regression as a one-layer network.}  When the
nonlinearity is the sigmoid $\sigma(x) = 1 / (1 + e^{-x})$ and the
output is interpreted as a probability, equation~\eqref{eq:nn_layer} is
exactly \emph{logistic regression} --- the classical statistical model
for binary classification.  This is not a coincidence: a neural network
is, layer for layer, a stack of generalised regressions, each producing
features for the next.

\paragraph{Multiple layers.}  Stacking layers gives a deep network:
\begin{equation}
    \zb^{(L)} \;=\; f_L\!\left( \Wb_L \, f_{L-1}\!\left( \Wb_{L-1}\,
    \cdots f_1\!\left( \Wb_1 \xb + \mathbf{b}_1 \right) \cdots
    + \mathbf{b}_{L-1} \right) + \mathbf{b}_L \right).
    \label{eq:nn_deep}
\end{equation}
Without the nonlinearities $f_i$, this composition would collapse into
a single linear transformation --- the product $\Wb_L \cdots \Wb_1$ ---
giving us no more expressive power than a one-layer network.  The
nonlinearity is what makes depth meaningful.  We met this argument
already in the Scattering Transform, where the absolute value played
the same role.

\marginnote{%
\textbf{The universal approximation theorem.}  A sufficiently wide
two-layer network with a non-polynomial nonlinearity can approximate
any continuous function on a compact domain to any desired accuracy.
This is a celebrated theoretical result, and it sounds powerful, but
it understates the role of depth.  A two-layer network may need
exponentially many neurons to do what a deeper network does with
polynomially many.  Depth is not just possible; it is efficient.}

\paragraph{How the weights are learned.}  This is the algorithmic heart
of modern machine learning.  Given a dataset of input-output pairs
$\{(\xb_n, \yb_n)\}$, we define a \emph{loss function} $\mathcal{L}$
that measures how badly the network's predictions match the targets.
We then adjust the weights to reduce the loss by \emph{gradient
descent}:
\begin{equation}
    \Wb_\ell \;\leftarrow\; \Wb_\ell \,-\, \eta\,
    \frac{\partial \mathcal{L}}{\partial \Wb_\ell},
\end{equation}
with a learning rate $\eta$.  Computing the gradients
$\partial \mathcal{L} / \partial \Wb_\ell$ for every layer is done by
\emph{backpropagation}, which is the multivariate chain rule applied
methodically from the output layer back to the input layer.  Part~II
of this book derives backpropagation in detail, implements it in C++,
and verifies it with numerical gradient checks.

\paragraph{Encoders and decoders.}  We can use depth-with-nonlinearity
to build two complementary structures:

\begin{itemize}
    \item An \emph{encoder} starts with a high-dimensional input and
    produces a progressively lower-dimensional representation,
    distilling the input to a compact latent code.  Each layer has
    fewer outputs than inputs.
    \item A \emph{decoder} reverses the process: starts with a small
    latent vector and produces a high-dimensional output.  Each layer
    has more outputs than inputs.
\end{itemize}

\noindent An \emph{autoencoder} chains the two: encoder, latent
bottleneck, decoder, trained so the output reconstructs the input.
The latent bottleneck forces the network to find a compact
representation that loses as little information as possible.  When the
encoder and decoder are both deep nonlinear networks, the autoencoder
is a learned analogue of PCA --- with the same goal (dimensionality
reduction) but able to capture curved manifolds rather than just flat
hyperplanes.

A \emph{variational autoencoder} extends the autoencoder to a
generative model: the bottleneck is interpreted as a probability
distribution, and we can sample new outputs by sampling from the
latent distribution and running them through the decoder.  Part~II
builds both, and the audio VAE chapter of Part~II uses one to morph
between sounds.

%------------------------------------------------------------------------
\section{Convolutional neural networks}
%------------------------------------------------------------------------

A convolutional neural network applies the inner product locally, like
the Scattering Transform of the previous chapter, but with the filters
learned from data rather than designed by hand.  The architectural
template is exactly the same:

\begin{itemize}
    \item Convolve the input with a bank of filters.
    \item Apply a nonlinearity.
    \item Possibly pool (spatial downsampling).
    \item Repeat for several layers.
    \item Finish with a few dense layers for classification.
\end{itemize}

\noindent What changes from the scattering picture is that the filters
are no longer wavelets chosen for analytical convenience; they are
free parameters trained jointly with the rest of the network.  Each
filter is a small inner product applied at every location in the input:
the same inner product, slid across the signal.

\paragraph{Why share weights.}  The defining feature of a CNN compared
to a fully-connected network is \emph{weight sharing}: the same filter
weights are used at every spatial position.  This has three
consequences.  It reduces the number of parameters dramatically (a
$3 \times 3$ filter on an image of any size has only 9 weights, not 9
weights per image location).  It builds in \emph{translation
equivariance}: shifting the input shifts the output.  And it encodes
a prior of locality: each output value depends on a small local
neighbourhood, capturing the fact that pixel relationships are
usually nearby.

These three are not optional features; they are the architectural prior
that makes CNNs work on images and audio while fully-connected networks
struggle.

\marginnote{%
\textbf{Empirically interpretable.}  Train a CNN on natural images and
the first-layer filters --- which are easy to visualise as small RGB
patches --- come out looking like oriented edge detectors and
colour-opponent blobs.  These are exactly the kinds of features that
classical signal-processing theory tells us \emph{should} be useful for
natural images.  Krizhevsky's AlexNet paper (2012) and Zeiler and
Fergus's 2014 follow-up made this picture famous; we reproduce it in
Part~II's CIFAR-10 example.}

\paragraph{The hierarchy.}  In a multi-layer CNN, early layers learn
low-level features (edges, colour patches, texture primitives), middle
layers learn combinations (corners, simple shapes, repeating textures),
and late layers learn task-specific concepts (object parts, faces,
musical-instrument families).  This hierarchical organisation emerges
naturally from gradient descent on a classification loss; nothing in
the training procedure explicitly asks for it.

The same hierarchical pattern appeared in the Scattering Transform:
first-layer wavelet coefficients capture frequency content, second-layer
coefficients capture modulations of that content.  The CNN learns the
same kind of hierarchy, but discovers what the relevant features
\emph{are} from the data rather than receiving them as wavelets.  This
is the gain from learning, and it is the gain that explains why CNNs
dominate image and audio tasks in practice.

%------------------------------------------------------------------------
\section{The unified picture}
%------------------------------------------------------------------------

We have now seen, across two chapters, a sequence of representations:

\begin{itemize}
    \item \textbf{DFT}: linear projection onto sinusoids.
    \item \textbf{DWT}: linear projection onto scaled-shifted wavelets.
    \item \textbf{Scattering}: multi-layer wavelet cascade with absolute
    value as nonlinearity.
    \item \textbf{PCA}: linear projection onto eigenvectors of the data
    covariance.
    \item \textbf{$K$-means}: piecewise-constant assignment to learned
    prototypes.
    \item \textbf{Matching pursuit}: sparse selection from a learned (or
    designed) overcomplete dictionary.
    \item \textbf{Fully-connected network}: stacked nonlinear projections
    with learned dense weights.
    \item \textbf{Convolutional network}: stacked nonlinear projections
    with learned local weights, shared across spatial position.
\end{itemize}

\noindent Every one of these fits the template $\Phi(\xb) = f(\mathbf{B}
\xb)$, possibly with several layers stacked, possibly with the
projection $\mathbf{B}\xb$ replaced by a convolution.  The two
axes of variation are: (1) where the reference vectors come from
(predefined or learned), and (2) whether locality is enforced
(global inner products versus local convolutions).  This gives a $2
\times 2$ table:

\begin{center}
\begin{tabular}{l|cc}
        & \emph{Global} (inner product) & \emph{Local} (convolution) \\ \hline
\emph{Designed basis} & DFT, PCA, $K$-means & DWT, Scattering \\
\emph{Learned basis}  & FC network, matching pursuit & CNN \\
\end{tabular}
\end{center}

The progression of this book is along the diagonal: from designed-global
(Chapter~\ref{ch:apriori}) through designed-local and learned-global
(this chapter) to learned-local (the CNN, and the conv-VAE we will
build in Part~II).  Each step adds capability; each step also adds
opacity, because learned representations are harder to interpret than
designed ones.

The next chapter closes Part~I by working a concrete problem ---
musical style transfer --- through each cell of this table.  We will
see what \emph{the same problem looks like} when posed in the waveform
(Fourier), in the spectrum (DFT), in the wavelet domain (DWT), via
sparse atoms (matching pursuit), and via a learned latent space
(autoencoder).  The exercise will make vivid how the choice of
representation determines what is musically possible.
